\section{Conclusão}

Nesse trabalho foi realizado o experimento idealizado por J. Franck e G. Hertz e publicado originalmente em
1914, com intuito de comprovar experimentalmente a teoria de Bohr sobre o modelo atômico. 

O experimento consiste em posicionar um ânodo entre um cátodo e uma placa coletora dentro de um tubo preenchido com 
vapor de mercúrio. 
O cátodo é aquecido, e assim elétrons são termicamente ejetados, formando uma nuvem eletrônica ao redor dele.
Esses elétrons são acelerados em direção ao ânodo por uma diferença de potencial. 
Alguns desses elétrons passam pelo ânodo e são coletados pela placa, desde que tenham energia suficiente para vencer
o potencial de retardo que há entre ela e o ânodo.

A corrente elétrica que vem da placa é medida, servindo como uma forma de contagem dos elétrons que atingem a placa.
Os resultados obtidos mostram que a corrente, como função do potencial entre o cátodo e o ânodo, apresenta uma série de 
picos regularmente espaçados. 

Isso implica que para valores bem definidos do potencial, ocorrem processos que reduzem a energia dos elétrons, sempre
pela mesma quantidade.

Esse fenômeno é muito bem explicado pelo modelo atômico de Bohr, onde os elétrons possuem órbitas estácionárias 
com energia bem definida, de modo que é necessária uma quantidade específica de energia para excitar o átomo. 
Embora esse modelo não esteja $100\%$ correto, e os elétrons não girem ao redor do núcleo como Bohr acreditava, é fato 
que os níveis de energia são quantizados, o que é demonstrado por esse experimento.

O que aconteceu foi que, ao aumentar o potencial acelerador para além de um certo valor crítico, 
os elétrons passam a poder transferir parte de sua energia para os átomos de Hg com que colidem, não sendo mais 
capazes de alcançar a placa coletora, e ocasionando uma queda na corrente.

Aumentando mais ainda o potencial de excitação, os elétrons voltam a chegar na placa, e a corrente passa a crescer.
Mas, para intervalos regulares de $V$, a corrente torna a cair, pois um mesmo elétron passa a poder transferir energia
mais de uma vez.
 
O produto do valor desse intervalo pela carga de um elétron é então a energia necessária para levar um átomo de Hg do seu estado fundamental para seu primeiro estado excitado. 
Pelos resultados obtidos, foi obtido o valor de $\qty{5,1}{\eV}$ para essa energia, o que equivale a um desvio de $5\%$ com relação ao valor aceito pela literatura de $\qty{4,9}{\eV}$.

Foi observado também que para $T$ suficientemente baixo, não ocorre nenhuma redução da corrente. 
Isso se deve ao fato de que se a pressão do gás for suficientemente pequena, colisões se tornam extremamente 
improváveis de ocorrer. 